\newcommand{\wineIntro}{Tonight, we are grateful for the unprecedented freedom
in our present day lives\textemdash{}far more freedom than our ancestors could ever have
dreamed was possible.}

\newcommand{\wineSummaryThisYear}{This year, we will hear some
strategies\textemdash{}adapted from David Cain's writings on
Raptitude\textemdash{}to feel deep gratitude for the freedoms we enjoy, that we
may better appreciate the importance of bringing freedom to the rest of the
world.}

\newcommand{\wineFirstCup}{
  \newReading With this first cup of wine, we consider how to feel more grateful
  for the freedom we enjoy 
  by noticing the experiences we have that happen because we are free.
  
  Many gratitude exercises ask you to think about what you have to be grateful
  for. In other words, you brainstorm reasons you ought to feel grateful,
  whether or not you do.

  These exercises are worthwhile on some level, but most of the time they don't
  create much of a real-time, visceral sense of gratitude. They just remind you
  of certain encouraging rote facts: on paper, your situation is pretty good;
  many parts of your life would be enviable to others; our ancestors certainly
  had things far worse as slaves in the Land of Egypt.

  \newReading As you might have noticed, simply making the case to ourselves
  that we have \textit{reasons} to feel grateful and celebrate our freedom
  doesn't necessarily make us feel that way.
  
  Gratitude, when we do genuinely feel it, arises from experiences we are
  \textit{currently having}, not from evaluating our lives in our heads. When
  you feel lonely, simply remembering that you have friends is a dull, nominal
  comfort compared to how wonderful it feels when one of those friends calls you
  out of the blue. Reflecting on the good fortune of having a fixed address is
  nice, but stepping inside your front door after a cold and rainy walk home is
  sublime. 
  
  The experience, not the idea, is what matters. So if you want to feel grateful
  for your freedom, move beyond the thinking exercises. Look for your good
  fortune not in some abstract assessment of your life situation, but in your
  experience right in this moment. What can you see, feel, hear, or sense, right
  here in the present, that's helpful, pleasant, or beautiful? Any interesting
  sensory experience or pleasant feeling will do: the cool breeze from a fan,
  the pet on your lap, the taste of this wine.
  
  \newReading That's the other important part: we don't need to reserve our
  gratitude for the big, fundamental conditions, % TODO (Eli 20220403): come up
  with better adjective for all-encompassing
  such as health, economic class, or loved ones. In every moment, regardless of
  your problems, you are enjoying freedom through all sorts of small, often
  haphazard pleasures: the smell of a meal wafting in from the kitchen, the
  softness of this chair, the alertness of your mind right now.
  
  Of course, all the abstract, big-picture life conditions made possible by your
  freedom have their own corresponding present-moment pleasures, and they are
  what matter. Consider the world of difference between trying to appreciate the
  notion that you aren't homeless, and appreciating the real-time experience of
  getting into bed in your own bedroom. That's where your good fortune truly
  resides\textemdash{}in your experiences, not your thoughts.
  
  With this first cup of wine, we remember to enjoy our present experiences to
  help us feel grateful for all the freedom that has been given to us.
}

\newcommand{\wineSecondCup}{
  \newReading With this second cup of wine, we learn a technique to more easily
  feel gratitude during everyday experiences; during a seemingly ordinary
  activity, describe it in a way that would be understandable to our ancestors
  in Egypt. 
  
  
  I wake up in a soft bed. I'm warm and safe, like every morning.
  As I rise and stretch, I notice I'm sore. Not from tending the fields though.
  I have no fields.
  Sometimes I forget that there's any field-tending going on at all. I buy all
  my food\textemdash{}I wouldn't know how to grow it or hunt it. I've never gone a day
  without food.
  
  My soreness is actually from my leisure time, not work. I spent yesterday
  sliding down a snow-covered slope with a board attached to my feet. After that
  I was pretty worn out, so I went to a friend's house, and watched a sporting
  event as it unfolded in Boston. I don't live in Boston, but my friend has a
  machine that lets us see what's happening there. I have one too. Almost
  everyone does.
  
  \newReading The sun won't rise for another hour, but I don't need to light a
  fire or candles. I have artificial ones, mounted on the ceiling. Hit a tiny
  switch and I can see everything, any time of day.
  
  I bathe while standing. The water comes out whatever temperature I like.
  
  I use a few machines in my kitchen to get my breakfast ready. It takes about
  five minutes.
  I don't know where it came from but I'd be surprised if it was from anywhere
  near here. As I eat, it occurs to me that my co-worker has a machine I might
  need to use at work today, so I want to make sure he brings it with him. He
  lives five miles from me, but that's no problem. I've got a device that lets
  him hear my voice from that distance. Wherever he is, I can talk to him.

  One minute later I've solved that problem, and forgotten about it.

  \newReading I put my dishes into another machine, which will clean them for me
  while I'm away at work.
  My destination is ten miles in the distance and I've got twenty-five minutes,
  which is plenty of time, because I won't be walking.

  I get into my most expensive machine. It's actually quite miraculous, but I
  forget that all the time. It allows me to sit in a comfortable chair, sealed
  from the elements, while it propels me at incredible speeds.

  After work, I hurtle home in the same manner, without really thinking about
  it.
  Then I turn on one of my favorite machines. It's about the size of a scrap of
  parchment.

  It has a glowing window inside it. A single page. But I only need one page
  because its contents change at my command. Sometimes there are words,
  sometimes pictures, sometimes both. Some pictures can move and talk.
  
  \newReading The stuff inside can be written by anyone in the world, even
  \textit{while it's in my hand}. There's more inside than I could ever read. It
  provides me with unbelievable advantages. Anything I don't know, I can find
  out in a few seconds.
  I can buy pretty much anything from where I'm sitting, and have it brought to
  my door.
  
  The sun went down hours ago, but with my artificial light I haven't noticed.
  When I'm able to summon the willpower, I close my favorite machine and go to
  bed.

  I'm warm and safe, like every night.
  
  We drink this second cup realizing that our ancestors would not have been able
  to even conceive of a life filled with as much freedom as we enjoy today. 
}

\newcommand{\wineThirdCup}{
  \newReading With the third cup of wine, we appreciate the tremendous benefits
  in our lives afforded by our freedom by imagining if they were to disappear.
  
  As a kid, whenever I stayed for supper at certain friends' houses, I wasn't
  sure what to do when they prayed.

  My family didn't say grace, but I knew a bit about the ritual from TV. I knew
  you were supposed to look down and say amen at the end, so I did.
  
  ``Grace isn't really religious,'' a formerly religious co-worker said to me
  years later.
  
  ``Yeah, you don't have to believe the food came from the Lord. All they're
  saying is, `Before we eat this food, let's remember that we could have no
  food.'{}''
  
  \newReading \textit{This might not be so}, in other words. This thing you have
  because you are free and not a slave in Egypt\textemdash{}this meal, this bed, these
  clothes, this friend\textemdash{}if it's possible to have it, it's possible not to have
  it. If you take a moment to imagine not having it, the good fortune of having
  it is no longer lost on you.
  
  Some nights, once I was comfortably in bed, I would imagine my bed
  disappearing around me and dropping me onto the bare floor. A moment later,
  the whole house would go, depositing me onto the cold, wet grass.
  
  I'd imagine this scenario in vivid detail. The dirt and dampness, the dead
  grass stuck to my arms. After a minute or so, the daydream would break, and
  the warmth and comfort of the covers\textemdash{}the happy reality of where I really
  was\textemdash{}would hit me like a truck.
  
  \newReading I'd do this on a smaller scale too, such as imagining my socks
  disappearing, just long enough to feel my clammy bare feet against the damp
  insoles of my boots. A moment later, I'd drop the fantasy, and discover anew
  just how much comfort a cheap pair of sport socks was adding to my life in
  that moment.
  
  I still do this. With socks, sweaters, walls. It's just one way to get your
  mind to collide with the essential truth behind any genuine experience of
  gratitude: this might not be so.  
  
  If you do the same exercise not with socks but with your loved ones, the
  effect is even more powerful. You'll need a moment in which you're with a
  loved one, and they're talking to someone else, doing a crossword, or
  otherwise not directly engaging with you.
  
  \newReading Whatever they're doing, quietly observe them doing it, paying
  attention to their unique mannerisms, voice, and presence.
  
  Now, try to see this moment as though it's actually in the past, and this
  person is gone now. You're remembering this one ordinary moment, in perfect
  detail, from that precious window of time when your lives still overlapped.
  
  It's important that it's an ordinary moment, because those are the kinds of
  moments that often seem neutral in value. Abundant and unremarkable. You might
  be tempted to spend it checking your phone.
  
  But if your mind can find that place, even for a just a second, where you see
  that \textit{this might not be so}, you'll break that sense that nothing
  special is happening.
  
  \newReading Then there will be a moment when you come back to reality. And the
  reality is that the precious, fleeting time when you got to be near this
  person\textemdash{}the time you'd do anything to revisit\textemdash{}is
  happening right now.
  
  The payoff isn't in the morbid fantasy about loss, it's the instant when you
  return from the vision and recognize what you really have, that you're living
  right in the middle of a fragile golden era. It's unique and beautiful and
  will one day be gone, and you are fantastically lucky to be here for it.
  
  With this third cup of wine, we feel gratitude for all of our freedoms and the
  joy they bring to our lives, by remembering that \textit{it might not be so}.
}

\newcommand{\wineFourthCup}{
  \newReading As our Seder draws near to the end, we take up our cups one last
  time as we think about how to feel grateful for all the freedom we enjoy even
  when faced with something unexpected.
  
  ``Radical Gratitude'' is about experimenting with being grateful for
  everything that happens to you. You're not telling yourself you
  \textit{should} feel grateful, only to invite or explore gratitude for what
  happens regardless of our initial feeling about it. The idea sounds
  ridiculous, and even hopeless, but in practice it's quite easy, and
  immediately rewarding. You just ask yourself, ``Can I be grateful for this
  too?'' In my short experience doing this, the brain has a way of coming up
  with good reasons why yes, you can.

  This practice reveals a lot about our short-sightedness. We have a rather
  ridiculous tendency to believe everything is either strictly good for us or
  bad for us, and that we can reliably determine which one it is, in the instant
  that thing happens.
  
  \newReading A few weeks ago, I went down to my storage locker and saw the door
  hanging open. My bike was gone. My initial feeling was the rush of violation
  and dirtiness that everyone feels when they see the mess left by a thief. They
  touched my stuff, and now some of it is at their place. But then I asked
  myself, ``Is there anything about this I could be grateful for?''
  
  The thought immediately put me into a totally different position, one where I
  didn't assume I should feel any particular way about it. The normal victim
  feelings gave way to something different. %a feeling of, ``Wow, I'm really
  glad I'm me.'' 
  I'm glad to know that the locker was so insecure before I put anything
  irreplaceable in there. I can afford a new bike. Also, I've never felt a
  desire to steal from people. Aren't I lucky that I don't know what it's like
  to enter a building illegally, and rifle through someone else's belongings? 
  
  \newReading Radical Gratitude is simply a way of challenging our initial
  feeling that a new development is wholly bad and that our moping and anger is
  justified, exploring instead what might also be good about it. Primarily, it
  does two things:

  It forces us out of the hypersensitive kind of autopilot we often operate
  under, which is based on a pretty grievous misconception: that events are
  isolated and are of two distinct types\textemdash{}good or bad\textemdash{}and
  that this goodness or badness is determined by how welcome it \textit{feels}
  when it happens.

  It also puts you into a helpful problem-solving state that often ends in
  gratitude for \textit{something} about what has just happened\textemdash{}the
  doors it opens, the things it teaches you, the future trouble it might spare
  you.

  \newReading Experimenting with this is also kind of fun. The more absurd it
  seems to be grateful, given the situation, the more interesting and fun it can
  be. Can I be grateful that my plans were canceled? Certainly. Can I be
  grateful I have a rash on my foot? Uh, let's find out.
  
  Again, being grateful for \textit{everything} is not an obligation; it's an
  option to experiment with. You aren't responsible for feeling any particular
  way after something happens.

  But that initial feeling\textemdash{}the horror of a locker door hanging
  open\textemdash{}doesn't have to be the final word on what's good and what's
  not. It's a question too complicated for our most reactive emotions to answer.

  With this fourth cup of wine, we consider how much easier it might be for us
  to find something to be grateful for in any event we're confronted with,
  compared to life for our ancestors in Egypt and those around the world still
  living in slavery today.
}
